% Part: free-logic
% Chapter: natural-deduction
% Section: common-rules

\documentclass[../../../include/open-logic-section]{subfiles}

\begin{document}

\olfileid{frl}{dis}{anf}

\olsection{Alternative Negative Free Logic}

To simplify matters, we only consider the Negative Free Logic above.

You should be aware that there is another way of developping Negative 
Free logic that allows $\eq[c][c]$ to be a logical truth even if $c$ doesn't 
exist. I'll discuss that variant briefly below, but you can ignore it.
The bottom line is: it's ugly and complicates things. 

On that alternative Negative Free Logic, we adopt the unrestricted
\Intro{\eq} rules as in Positive Free Logic. So $\eq[c][c]$ is a
logical truth; it holds even if $c$ is empty (doesn't denote
anything). The logic is still a \emph{negative} free logic because it
retains \Log{NFL}: \emph{atomic} claims like $\Obj{F}a$ cannot be true
if $a$ is empty. But we drop the analogue rule for $\eq$,
\Log{NFL\eq}. Thus in that version of Negative Free Logic:
\begin{itemize}
\item $\Obj{F}a\land\lnot\lfrexists a$ is inconsistent, as in our
original version of Negative Free Logic.
\item $\eq[a][a]\land\lfrexists a$ is not inconsistent.
\end{itemize}
So far this might seem nice. Unfortunately, ugly complications arise.
To make $\eq[a][a]$ a logic truth, we have to change the NFL semantics
for $\eq$. We face a choice: keeping the negative semantics in which
empty names don't refer to any object, or adopting the postive semantics 
on which empty names refer to 'non-existing' objects (the `outer' domain). 
Both options have serious downsides. 

If we keep the negative semantics, to make $\eq[c][c]$ a logical truth
we say that it's always true even if $c$ is empty. But when $c,d$ are 
both empty, what do we say about $\eq[c][d]$? The negative semantics 
has only two options: all of these are true, or all false. If all 
false, not only that feels artificial, but that requires a new ugly
rule: 
\begin{prooftree}
    \AxiomC{$\eq[c][d]$}
    \RightLabel{where $c$ and $d$ are \emph{distinct} constants}
    \UnaryInfC{$\lfrexists c$}    
\end{prooftree}
If all true, that might still feel artificial, and we also have to 
accept a new unappealing rule:
\begin{prooftree}
    \AxiomC{$\lnot\lfrexists c$}
        \AxiomC{$\lnot\lfrexists d$}
    \RightLabel{??}
    \BinaryInfC{$\eq[c][d]$}    
\end{prooftree}
Bottom line: with the negative semantics, saying that $\eq[c][c]$ 
is false when $c$ is empty might sound counterintuitive, but the 
other options aren't great either.

Let's consider the positive semantics instead. On that option we could
make $\eq[c][c]$ a logical truth even when $c$ is empty, \emph{and} 
saying that when $c,d$ are empty, $\eq[c][d]$ may or may not be true.
We would avoid the two arbitrary-seeming options above. This option 
is Positive Free Logic (rules and semantics) with the added
the signature law of Negative Free Logic, $\Log{NFL}$. That solution 
seems to have the worst of both worlds. Because it has the \Log{NFL} rule, 
its logic doesn't obey the nice principle of uniform substitution (unlike 
propositoinal logic, modal logic, first order logic and Positive Free Logic).
Because it uses the semantics of Positive Free Logic, it has a semantics
that assigns objects to empty names, hence a semantics that we cannot 
take at face value (unless we adopt a strange Meinongian metaphysics 
of things that are but don't exist). It's not clear what the motivation
would be for adopting that variant of Negative Free Logic over the 
simpler Positive Free Logic. 

\end{document}
